%%% =====================================================================
%%% figures1.tex  --  TikZ figures, set 1:
%%%   parton model, DIS, forward Compton amplitude, omega plane, OPE,
%%%   operator definition of PDFs, one-loop PDF graphs, DGLAP, splitting
%%%   vertices, O(alpha_s) DIS graphs, factorization schematic.
%%% (research/tikz1; local style prefix "t1")
%%%
%%% Requirements: preamble_fact.tex, or equivalently
%%%   \usepackage{tikz}
%%%   \usetikzlibrary{decorations.pathmorphing,decorations.markings,arrows.meta,
%%%                   calc,positioning,shapes.geometric,patterns,backgrounds,fit}
%%%   %%% BEGIN STYLES
% =====================================================================
%  Plain-TikZ Feynman-diagram styles (no tikz-feynman).
%  Load in the preamble, after \usepackage{tikz}:
%    \usetikzlibrary{decorations.pathmorphing,decorations.markings,arrows.meta,calc,positioning}
%  Contents of this block (all low-level names start with \fd):
%   (i)  helper registers/macros and three decorations
%          "fd gluon"      coil for open paths,
%          "fd gluon loop" coil for closed paths (full circles),
%          "fd photon"     sine wave;
%        they rescale the nominal segment length so that an integer
%        number of coils / half-waves fits the path exactly (no straight
%        left-overs, evenly spaced coils, also on arcs);
%   (ii) a single \tikzset with the user-level styles:
%          fd             picture style (line width, caps, label font)
%          quark[=pos]    solid line, centred Stealth arrow at pos (default 0.5)
%          gluon          coiled line;   gluon loop[=m]: closed coil, number
%                         of coils a multiple of m (use for full circles)
%          photon         wavy line
%          ghost[=pos]    dotted line with centred arrow
%          mom={..}, mom'={..}   momentum arrow parallel to a path, on its
%                         left / right:  \path[mom={$p$}] (A) -- (B);
%                         keys: mom pos, mom length, mom sep
%          vertex         vertex dot (node style):  \node[vertex] at (A) {};
%          blob           grey filled region:  \draw[blob] (c) circle[radius=5mm];
%          bound state    double line (e.g. a bound state)
% =====================================================================
\newdimen\fdAA \newdimen\fdTW \newdimen\fdHW
\def\fdloopmult{1}
\def\fdpt#1#2#3{\pgfqpoint{\dimexpr#1\fdAA+#3\fdTW\relax}{#2\pgfdecorationsegmentamplitude}}
\def\fdcoilfull{%
  \pgfpathcurveto{\fdpt{0}{0.555}{1}}{\fdpt{0.445}{1}{2}}{\fdpt{1}{1}{3}}%
  \pgfpathcurveto{\fdpt{1.555}{1}{4}}{\fdpt{2}{0.555}{5}}{\fdpt{2}{0}{6}}%
  \pgfpathcurveto{\fdpt{2}{-0.555}{7}}{\fdpt{1.555}{-1}{8}}{\fdpt{1}{-1}{9}}%
  \pgfpathcurveto{\fdpt{0.445}{-1}{10}}{\fdpt{0}{-0.555}{11}}{\fdpt{0}{0}{12}}}
\def\fdcoilhalf{%
  \pgfpathcurveto{\fdpt{0}{0.555}{1}}{\fdpt{0.445}{1}{2}}{\fdpt{1}{1}{3}}%
  \pgfpathcurveto{\fdpt{1.555}{1}{4}}{\fdpt{2}{0.555}{5}}{\fdpt{2}{0}{6}}}
\def\fdcoilopeninit{%
  \fdAA=\pgfdecorationsegmentaspect\pgfdecorationsegmentamplitude\relax
  \pgfmathsetmacro\fdN{max(1,round((\pgfdecoratedpathlength-2*\fdAA)/\pgfdecorationsegmentlength-0.5))}%
  \pgfmathsetlength\pgfdecorationsegmentlength{(\pgfdecoratedpathlength-2*\fdAA-0.3pt)/(\fdN+0.5)}%
  \fdTW=0.0833333\pgfdecorationsegmentlength\relax}%
\def\fdcoilclosedinit{%
  \fdAA=\pgfdecorationsegmentaspect\pgfdecorationsegmentamplitude\relax
  \pgfmathsetmacro\fdN{\fdloopmult*max(1,round(\pgfdecoratedpathlength/(\fdloopmult*\pgfdecorationsegmentlength)))}%
  \pgfmathsetlength\pgfdecorationsegmentlength{(\pgfdecoratedpathlength-0.1pt)/\fdN}%
  \fdTW=0.0833333\pgfdecorationsegmentlength\relax}%
\def\fdphotoninit{%
  \pgfmathsetmacro\fdN{max(1,round(2*\pgfdecoratedpathlength/\pgfdecorationsegmentlength))}%
  \pgfmathsetlength\fdHW{(\pgfdecoratedpathlength-0.1pt)/\fdN}}%
\pgfdeclaredecoration{fd gluon}{init}{%
  \state{init}[width=+0pt,next state=coil,persistent precomputation=\fdcoilopeninit]{}%
  \state{coil}[switch if less than=1.5\pgfdecorationsegmentlength+2\fdAA to last,
               width=+\pgfdecorationsegmentlength]{\fdcoilfull}%
  \state{last}[width=+\dimexpr.5\pgfdecorationsegmentlength+2\fdAA\relax,next state=final]{\fdcoilhalf}%
  \state{final}{\pgfpathlineto{\pgfpointdecoratedpathlast}}}
\pgfdeclaredecoration{fd gluon loop}{init}{%
  \state{init}[width=+0pt,next state=coil,persistent precomputation=\fdcoilclosedinit]{}%
  \state{coil}[switch if less than=+\pgfdecorationsegmentlength to final,
               width=+\pgfdecorationsegmentlength]{\fdcoilfull}%
  \state{final}{\pgfpathlineto{\pgfpointdecoratedpathlast}}}
\pgfdeclaredecoration{fd photon}{init}{%
  \state{init}[width=+0pt,next state=up,persistent precomputation=\fdphotoninit]{}%
  \state{up}[switch if less than=+\fdHW to final,width=+\fdHW,next state=down]{%
    \pgfpathsine{\pgfqpoint{.5\fdHW}{\pgfdecorationsegmentamplitude}}%
    \pgfpathcosine{\pgfqpoint{.5\fdHW}{-\pgfdecorationsegmentamplitude}}}%
  \state{down}[switch if less than=+\fdHW to final,width=+\fdHW,next state=up]{%
    \pgfpathsine{\pgfqpoint{.5\fdHW}{-\pgfdecorationsegmentamplitude}}%
    \pgfpathcosine{\pgfqpoint{.5\fdHW}{\pgfdecorationsegmentamplitude}}}%
  \state{final}{\pgfpathlineto{\pgfpointdecoratedpathlast}}}
\tikzset{
  /pgf/decoration/coil multiple/.store in=\fdloopmult,
  fd/.style={line width=0.8pt,line cap=round,line join=round,font=\small},
  fd arrow/.style={postaction={decorate,solid,line width=0.8pt,decoration={markings,
      mark=at position #1 with {\arrow[xshift=2.9pt]{Stealth[length=5.5pt,width=4.5pt]}}}}},
  quark/.style={draw,solid,line width=0.8pt,fd arrow=#1},
  quark/.default=0.5,
  plain/.style={draw,solid,line width=0.8pt},
  gluon/.style={draw,solid,line width=0.75pt,decorate,
      decoration={fd gluon,amplitude=3.2pt,segment length=4.6pt,aspect=0.7}},
  gluon loop/.style={draw,solid,line width=0.75pt,decorate,
      decoration={fd gluon loop,amplitude=3.2pt,segment length=4.6pt,aspect=0.7,coil multiple=#1}},
  gluon loop/.default=1,
  photon/.style={draw,solid,line width=0.8pt,decorate,
      decoration={fd photon,amplitude=2pt,segment length=6.6pt}},
  ghost/.style={draw,line width=1.3pt,line cap=round,dash pattern=on 0pt off 2.6pt,fd arrow=#1},
  ghost/.default=0.5,
  mom pos/.store in=\fdmompos,mom pos=0.5,
  mom length/.store in=\fdmomlen,mom length=0.55cm,
  mom sep/.store in=\fdmomsep,mom sep=0.25cm,
  mom arrow/.style={draw,solid,line width=0.5pt,-{Stealth[length=3.8pt,width=3.2pt]}},
  mom/.style={decorate,decoration={markings,mark=at position \fdmompos with {%
      \draw[mom arrow] (-0.5*\fdmomlen,\fdmomsep) -- node[auto,inner sep=1.5pt]{#1} (0.5*\fdmomlen,\fdmomsep);}}},
  mom'/.style={decorate,decoration={markings,mark=at position \fdmompos with {%
      \draw[mom arrow] (-0.5*\fdmomlen,-\fdmomsep) -- node[auto,swap,inner sep=1.5pt]{#1} (0.5*\fdmomlen,-\fdmomsep);}}},
  vertex/.style={circle,fill,draw=none,inner sep=0pt,minimum size=3.8pt},
  blob/.style={draw,solid,fill=gray!30,line width=0.8pt},
  bound state/.style={draw,solid,double,double distance=1.6pt,line width=0.6pt},
}
%%% END STYLES
%%% BEGIN STYLES-EXT
% =====================================================================
%  Extension of the plain-TikZ Feynman-diagram library (fd_styles.tex)
%  for the factorization monograph.  Requires fd_styles.tex loaded first.
%  Extra libraries: shapes.geometric, patterns, backgrounds, fit.
%  New user-level styles:
%    wilson[=pos]   Wilson (eikonal) line: double line with centred arrow
%    wilson plain   Wilson line without arrow
%    cut            final-state cut: dashed red vertical line
%    hard blob      hard subgraph (grey);   jet blob (A side, blue);
%    jetB blob      collinear subgraph of hadron B (orange);
%    soft blob      soft subgraph (green);  hadron blob (target, light grey)
%    glauber        Glauber gluon (dark-red coil)
%    opvertex       operator insertion vertex (circled cross), node style
%    scalar         dashed line (e.g. scalar/longitudinal polarization)
%    longit         gluon carrying k^mu (scalar) polarization: blue coil;
%                   optionally mark it with \longitmark{(A)}{(B)} (dot at midpoint)
%    lcaxis         thin arrow for light-cone / spacetime axes
%    region         dashed rounded rectangle marking a momentum region
% Colours are kept light so that the figures print well in greyscale.
% =====================================================================
\tikzset{
  wilson/.style={draw,solid,double,double distance=1.6pt,line width=0.55pt,fd arrow=#1},
  wilson/.default=0.5,
  wilson plain/.style={draw,solid,double,double distance=1.6pt,line width=0.55pt},
  cut/.style={draw,dashed,red!70!black,line width=0.9pt,dash pattern=on 4pt off 2.5pt},
  hard blob/.style={draw,solid,fill=gray!35,line width=0.8pt},
  jet blob/.style={draw,solid,fill=blue!12,line width=0.8pt},
  jetB blob/.style={draw,solid,fill=orange!18,line width=0.8pt},
  soft blob/.style={draw,solid,fill=green!15,line width=0.8pt},
  hadron blob/.style={draw,solid,fill=gray!15,line width=0.8pt},
  glauber/.style={draw,solid,line width=0.75pt,color=red!60!black,decorate,
      decoration={fd gluon,amplitude=3.2pt,segment length=4.6pt,aspect=0.7}},
  opvertex/.style={circle,draw,solid,fill=white,line width=0.8pt,inner sep=0pt,minimum size=7pt,
      path picture={\draw[line width=0.6pt] (path picture bounding box.north west)--(path picture bounding box.south east)
                    (path picture bounding box.south west)--(path picture bounding box.north east);}},
  scalar/.style={draw,solid,line width=0.8pt,dash pattern=on 3pt off 2pt},
  longit/.style={draw,solid,line width=0.75pt,color=blue!60!black,decorate,
      decoration={fd gluon,amplitude=3.2pt,segment length=4.6pt,aspect=0.7}},
  lcaxis/.style={draw,solid,line width=0.5pt,-{Stealth[length=4pt,width=3.4pt]}},
  region/.style={draw,dashed,rounded corners=3pt,line width=0.6pt,gray!70!black},
}
\newcommand{\longitmark}[2]{\fill[blue!60!black] ($#1!0.5!#2$) circle[radius=1.7pt];}
%%% END STYLES-EXT

%%% plus subcaption, cleveref (\cref) and the preamble macros
%%%   \ii \dd \ee \Disc \Real \Imag \as \order.
%%% Every extra style is set with \tikzset INSIDE the figure environment
%%% (hence local to that figure) and carries the prefix "t1":
%%%   t1 hadron, t1 parton, t1 dot, t1 boost, t1 bcut, t1 contour, t1 carrow.
%%% Helper macros \toneW, \toneDIS are \newcommand'ed inside their figure
%%% environments (local to the figure group); the scratch macros \toneR,
%%% \toned, \tonea, \tonex, \toneX are defined inside tikzpicture groups.
%%%
%%% Drawing conventions used throughout this set
%%%   * solid line with arrow  = quark; the arrow is the fermion-number flow;
%%%   * thin separate arrow with label (styles mom/mom') = momentum flow;
%%%   * thick line (t1 hadron) = hadron; double line = Wilson (eikonal) line;
%%%   * red dashed vertical line = final-state cut (amplitude on the left,
%%%     complex-conjugate amplitude on the right);
%%%   * zig-zag line on a real axis = branch cut of an analytic function;
%%%   * where two lines cross without a vertex, the one drawn on top is
%%%     separated by a short white gap in the other (fig:forward-compton b).
%%% =====================================================================
